\documentclass[10pt,a4paper]{amsart}
\usepackage[margin=19mm]{geometry}
\usepackage{amsmath,amssymb,mathtools}
\usepackage[scr=rsfs,cal=euler]{mathalfa}
\usepackage{xfrac}
\usepackage[unicode,hidelinks,bookmarks=false]{hyperref}
\newcommand{\ie}{i.e.\ }
\newcommand{\cf}{cf.\ }
\newcommand{\resp}{respectively\ }
\newcommand{\Subsection}{\subsection}
\providecommand{\nobreakdash}{\nobreak\hbox{-}}
\setlength{\emergencystretch}{2em}
\multlinegap0pt
\usepackage{amssymb}
\usepackage{mathtools}
\usepackage{xcolor}
\usepackage{bm}
\usepackage{algorithm}\usepackage{algpseudocode}
\usepackage{cancel}
\usepackage{xfrac}
\usepackage[shortlabels]{enumitem}
\newtheorem{lemma}{Lemma}
\newcommand\bfF{{\mathbf F}}
\newcommand\bfS{{\mathbf S}}
\newcommand\bfU{{\mathbf U}}
\newcommand\bfV{{\mathbf V}}
\newcommand\bfY{{\mathbf Y}}
\newcommand\bfZ{{\mathbf Z}}
\title{High-order robust basis-update \& Galerkin integrators for dynamical low-rank approximation}
\author{Cory D. Hauck, Jonas Kusch, Steffen Schotth\"ofer}
\date{Source: arXiv:2608.27749v1 (CC BY 4.0)}
\begin{document}
\maketitle

\noindent\emph{Source and licence.} High-order robust basis-update \& Galerkin integrators for dynamical low-rank approximation. Version arXiv:2608.27749v1 (CC BY 4.0), distributed by arXiv under the Creative Commons Attribution 4.0 International licence, http://creativecommons.org/licenses/by/4.0/, as recorded in the arXiv licence metadata for this exact version and shown on the abstract page https://arxiv.org/abs/2608.27749v1. Full licence notice: \texttt{licenses/arxiv-2608.27749-CC-BY-4.0.txt}.\par\smallskip

\noindent\textit{Editorial scope.} Counted unit: the article's lemma le:approxU (source0.tex lines 345--347). The article's complete original proof of it is delivered with it (source0.tex lines 348--366, one \texttt{proof} environment of 19 lines). No mathematics is written, repaired or re-derived: the statement and the proof are the article's own lines, in the article's own notation, and no retained line is substituted or re-flowed.  Retained uncounted as the source-local context the proof needs: lines 339--342.  Nothing was omitted from any retained span, so the package carries no omission record.  No retained line contains a citation, so the wrapper carries no bibliography and no bibliography entry.  No retained line contains a figure environment, an image-inclusion command or drawing code, so no image is delivered and the package declares no asset dependency.  Editorial formatting only, in addition: an AMS \texttt{amsart} preamble that restates the article's own declarations and macros used by the retained lines, namely the environments \texttt{lemma} on separate counters, together with the standard packages those lines need.  The retained environments print in this document as Lemma 1 in order of appearance, whereas the article prints the same items with the numbering of its own full document.  The complete native source of the article as distributed in the arXiv e-print for this exact version is bundled unchanged as \texttt{sources/208736/source0.tex}.
\begin{align}
    \widehat\bfU_{(j)} :=\,& \text{orth}([\bfU_0, \,\bfF_{(1)}\bfV_{(1)}, \,\bfF_{(2)}\bfV_{(2)}, \,\cdots \,, \,\bfF_{(j)}\bfV_{(j)}])\,,\label{eq:augBasisGeneralU} \\ 
    \widehat\bfV_{(j)} :=\,& \text{orth}([\bfV_0, \,\bfF_{(1)}^{\top}\bfU_{(1)}, \,\bfF_{(2)}^{\top}\bfU_{(2)}, \,\cdots\,, \,\bfF_{(j)}^{\top}\bfU_{(j)}])\,. \label{eq:augBasisGeneralV}
\end{align}

\begin{lemma}[Basis information]\label{le:approxU}
    $\bfU_{(\ell)}$ is spanned by $\widehat \bfU_{(j-1)}$ and $\bfV_{(\ell)}$ is spanned by $\widehat \bfV_{(j-1)}$ for $\ell \leq j$.
\end{lemma}

\begin{proof}
We use induction to show this result holds for general $j\in\mathbb{N}$. The statement is true for $j=1$ since with $\widehat \bfU_{(0)} \equiv \bfU_{0}$ we have
\begin{align*}
    \bfU_{(1)} = \bfU_{0} \in \text{span}(\widehat \bfU_{(0)})\,.
\end{align*}
Let the statement hold for an arbitrary but fixed $j$ (IH). Then, since
\begin{align}
    \bfY_{(j+1)}:= \bfU_{(j+1)}\bfS_{(j+1)}\bfV_{(j+1)}^{\top} := [\![\bfY_0 + h\sum_{\ell = 1}^{j}a_{j,\ell}\bar \bfZ_{\ell}]\!]\,,
\end{align}
we know that 
\begin{align*}
    \bfU_{(j+1)} \in\,& \text{span}\{ \bfU_0, \bar \bfZ_1, \cdots,  \bar \bfZ_j\}\,\\
    =\,& \text{span}\{ \bfU_0, \bfF_{(1)}\bfV_{(1)},  \bfU_{(2)}, \bfF_{(2)}\bfV_{(2)}, \cdots,  \bfU_{(j)}, \bfF_{(j)}\bfV_{(j)}\}\,\\
    \stackrel{\mathrm{(IH)}}{=}\,& \text{span}\{ \widehat \bfU_{(j-1)}, \bfF_{(1)}\bfV_{(1)}, \cdots,  \bfF_{(j)}\bfV_{(j)}\}\\
    \stackrel{\mathrm{\eqref{eq:augBasisGeneralU}}}{=}\,& \text{span}\{ \widehat \bfU_{(j-1)}, \bfF_{(j)}\bfV_{(j)}\}\\
    \stackrel{\mathrm{\eqref{eq:augBasisGeneralU}}}{=}\,& \text{span}\{ \widehat \bfU_{(j)}\}\,.
\end{align*}
Showing that $\bfU_{(\ell)}\in \text{span}\{ \widehat \bfU_{(j)}\}$ for $\ell\leq j$ directly follows from the discussion above. The proof for $\bfV_{(\ell)}$ follows analogously.\qed
\end{proof}

\end{document}
